\section{Introduction}\label{main:introduction}

Categoricity asks whether cardinality determines a model up to isomorphism.
For abstract elementary classes, the eventual categoricity conjecture predicts
that uniqueness in one sufficiently large cardinal forces uniqueness in every
cardinal above a uniform threshold. We prove this prediction in ZFC, with the
threshold depending only on a bound for the L\"owenheim--Skolem number.

\subsection{The statement}

We specify the axioms, including the union clause used throughout the proof.
All structures have sets as their universes, and all languages are set-sized
and finitary. We write $|M|$ for the universe of a structure and $\|M\|$ for
its cardinality.

\begin{definition}[Abstract elementary class]\label{main:aec}
An \emph{abstract elementary class} (AEC) is a pair $K=(\mathcal C,\le_K)$,
where $\mathcal C$ is a class of structures in a fixed language $\tau$ and
$\le_K$ is a partial order satisfying the following conditions.
\begin{enumerate}[label=\textup{(\arabic*)}]
\item The class is closed under isomorphism. If $M\le_K N$, then $M$ is a
$\tau$-substructure of $N$. If $f:N\cong N'$ and $M\le_K N$, then
$f[M]\le_K N'$.
\item \emph{Coherence:} if $M_0$ is a $\tau$-substructure of $M_1$, and
$M_0\le_K M_2$ and $M_1\le_K M_2$, then $M_0\le_K M_1$.
\item For every nonzero ordinal $\delta$ and increasing chain
$(M_i:i<\delta)$ under $\le_K$, its union $U$ belongs to $\mathcal C$ and
$M_i\le_K U$ for all $i<\delta$. If also $M_i\le_K N$ for every $i$, then
$U\le_K N$.
\item There is an infinite cardinal $\ell\ge |\tau|+\aleph_0$ such that,
for every $M\in\mathcal C$ and $A\subseteq |M|$, there is $N\le_K M$ with
$A\subseteq |N|$ and $\|N\|\le |A|+\ell$.
\end{enumerate}
The least possible $\ell$ is denoted $\LS(K)$. A \emph{strong map} is an
isomorphism onto a strong substructure of its codomain. The class is
\emph{categorical in $\kappa$} if it has a model of cardinality $\kappa$ and
all its models of that cardinality are isomorphic.
\end{definition}

The chain axioms equivalently give unions and the common-upper-bound clause
for nonempty set-indexed directed systems of strong inclusions. We also use
the corresponding direct limits of coherent strong maps, obtained by
renaming their injective images. No intersection axiom is included.

\begin{theorem}[Uniform eventual categoricity]\label{main:uniform}
For every infinite cardinal $\lambda$ there is a cardinal
$\mu(\lambda)\ge\lambda$ such that the following holds for every AEC $K$
with $\LS(K)\le\lambda$:
\[
  \bigl(\exists\kappa\ge\mu(\lambda)\quad K\text{ is categorical in }\kappa\bigr)
  \quad\Longrightarrow\quad
  \bigl(\forall\kappa'\ge\mu(\lambda)\quad
       K\text{ is categorical in }\kappa'\bigr).
\]
The same threshold is used in the premise and the conclusion.
\end{theorem}

In particular, Theorem~\ref{main:uniform} resolves Shelah's eventual
categoricity conjecture positively. No amalgamation, joint embedding,
tameness, no-maximal-models, or arbitrarily-large-models hypothesis is
assumed. The starting categoricity cardinal may be a successor, a regular
limit, or a singular cardinal. The theorem asserts the existence of a
threshold, without prescribing a numerical bound for it.

On the resulting categorical tail, every model strongly embeds into every
model of at least its cardinality. Indeed, the L\"owenheim--Skolem axiom
supplies a strong submodel of the smaller cardinality, and categoricity
identifies it with the given model. Taking a larger target gives proper
strong extensions and a common target for any two models. Thus the tail
has joint embedding and no maximal models.

It is useful to separate the uniform quantifier from the main structural
argument.

\begin{theorem}[The qualitative form]\label{main:qualitative}
An AEC categorical in unboundedly many cardinals is categorical on a tail of
the cardinals.
\end{theorem}

Here ``unbounded'' means unbounded in the class of all cardinals, and a tail
contains every cardinal above some fixed cardinal. The set-coding argument
in Proposition~\ref{found:uniform-reduction} proves the equivalence of the
two theorems. Once that reduction is made, arbitrarily large models are
available. The order construction in Theorem~\ref{found:order-template}
then supplies actual models at every sufficiently large cardinal; subsequent
categoricity arguments concern both existence and uniqueness.

The uniform theorem applies in particular to sentences of
$L_{\omega_1,\omega}$ in countable languages, where countable conjunctions
and disjunctions are allowed but quantifier strings are finite.
For such a sentence $\varphi$, choose a countable fragment containing it
and order its models by elementarity for that fragment. This is an AEC
with L\"owenheim--Skolem number $\aleph_0$
\cite[Example~2.0.2]{BoneyVasey2017Survey}. The models and their
isomorphisms are unchanged, so the same $\mu(\aleph_0)$ works for all
these sentences. The prescribed $\beth_{\omega_1}$ threshold in the
classical infinitary formulation remains a separate question
\cite[Conjecture~1.1]{BoneyVasey2015}.

\subsection{History and scope}

Morley's categoricity theorem gives the countable first-order prototype:
a countable theory categorical in one uncountable cardinal is categorical
in every uncountable cardinal \cite{Morley1965}.  Its structural approach
connects uniqueness of models with stability and dimension.  The geometry
of strongly minimal sets, developed by Baldwin and Lachlan
\cite[Sections~1--2]{BaldwinLachlan1971}, provides a central antecedent for
classification by bases.  In an AEC, elementary compactness and the
first-order description of types are not part of the axioms.  The
corresponding amalgamation, independence, and geometry therefore require
separate constructions.

Shelah's eventual categoricity conjecture asks for an eventual dichotomy
between categoricity and noncategoricity; see
\cite[Conjecture~N.4.2]{Shelah2009} and \cite{BoneyVasey2015}.
The equivalent formulation using one threshold depending on an LS bound
is discussed in \cite[Remark~1.3]{VaseyUniversalI}.  Our set-coding
argument makes that reduction explicit: it treats bounded categoricity
spectra as well as classes categorical at unboundedly many cardinals.
An explicit numerical threshold is a further question, separate from the
existential threshold in Theorem~\ref{main:uniform}; compare
\cite[Conjecture~N.4.3(1)(b)]{Shelah2009}.
Under CH, the companion \cite[Theorem~1.1]{OpenAI2026CHThreshold}
constructs an AEC $K$ with $\LS(K)=\aleph_0$ that is categorical on a
higher tail but not at $\beth_{(2^{\aleph_0})^+}=\beth_{\omega_2}$,
ruling out this prescribed endpoint for arbitrary AECs.

Substantial special cases have been established. Makkai and Shelah
proved transfer from sufficiently large successor categoricity for
$L_{\kappa,\omega}$-theories when $\kappa$ is strongly compact
\cite[Conclusion~5.1]{MakkaiShelah1990}. Grossberg and VanDieren
obtained upward categoricity transfer from one sufficiently large successor
cardinal for tame AECs with amalgamation, joint embedding, and arbitrarily
large models \cite[Theorem~5.2]{GrossbergVanDieren2006}.
Boney subsequently derived tameness from strongly compact cardinals and,
combining it with the Grossberg--VanDieren transfer, established eventual
categoricity for arbitrary AECs with a successor starting cardinal under
a proper class of strongly compact cardinals \cite{Boney2014}.
These results explain the importance of removing restrictions on the
starting cardinal as well as structural hypotheses on the class.

Universal classes admit
a ZFC transfer without assuming amalgamation
\cite[Theorem~7.3]{VaseyUniversalII}.  Tameness, amalgamation, and primes
give another transfer \cite[Theorem~3.8]{VaseyTamePrimes}.
Shelah and Vasey prove a ZFC categoricity theorem for excellent AECs and
obtain broader results using strongly compact cardinals, or the weak
generalized continuum hypothesis (WGCH), that is,
$2^\kappa<2^{\kappa^+}$ for every infinite cardinal $\kappa$, together
with amalgamation
\cite[Theorem~14.2 and Corollaries~14.7, 14.10]{ShelahVasey2024}.
Under WGCH, amalgamation, and arbitrarily large models, Vasey also gives
a classification of categoricity spectra
\cite[Corollary~9.7]{VaseySpectrum}.  These hypotheses supply substantial
structure. Here the amalgamation, bounded-tuple calculus, and geometry
needed for the working schemes are constructed at selected levels from
the AEC axioms and unbounded categoricity.

The proof uses several established methods in forms that are proved
locally.  Order presentations originate in the work of Ehrenfeucht and
Mostowski \cite{EhrenfeuchtMostowski1956}; their AEC formulation and
functorial maps are developed in
\cite[IV.0.8--IV.0.9 and IV.0.12]{Shelah2009}.  We retain a finite-support
presentation throughout the proof.  Its extraction uses classical
partition-calculus methods \cite{ErdosRado1956}, while the later
isolation argument uses Fodor's pressing-down method \cite{Fodor1956}.
Stationary independence and local character are central to Shelah's good
frames \cite[II.2.1]{Shelah2009}.  Sections~\ref{calc:section-calculus}
and~\ref{iso:section-isolation} establish the bounded-tuple calculus,
relative isolation, and dominated extensions required for the systems
used here.

A version-sensitive issue matters for the order argument.  Boney and
Vasey \cite[Fact~1.4 and the discussion of Claim~1.5, pp.~3--4]{BoneyVasey2015}
identify a gap in the unrestricted syntactic-characterization claim of
Shelah's Chapter~IV and withdraw an earlier attempted repair.  Their
valid variants with exponentiation or cofinality assumptions do not
cover all the starting cardinals of Theorem~\ref{main:uniform}.  We prove
the required reversal of strong maps directly in
Theorem~\ref{found:reversal}.

Finite systems and dimension-raising arguments have a long history in
nonelementary categoricity.  Shelah's existence and uniqueness methods
\cite[Sections~3--5]{Shelah1983B} are antecedents of the passage from
one additional dimension to a larger cardinal.  The modern treatment in
\cite[Definition~8.16 and Lemma~11.6]{ShelahVasey2024} distinguishes
amalgamated completion from strong uniqueness onto a prescribed target.
Both properties are needed below and are proved separately.  Our final
cardinal induction uses exact boundary maps for systems whose basis
pieces may have different cardinalities.

Esp\'indola has announced eventual categoricity for arbitrary AECs and
further accessible-category results \cite{Espindola2019}, and subsequently
an explicit-threshold categoricity theorem \cite[Theorem~4.1]{Espindola2023}.
These are separate public claims.  The argument here constructs its
working-scheme calculus, isolation, marked and geometric boundary lifting,
and common-level closure directly.  None of these interfaces is imported
from those preprints.

\subsection{Proof overview}

By the set-coding reduction in Section~\ref{found:section-reduction-orders},
it suffices to fix one AEC $\K$ categorical in unboundedly many cardinals
and prove categoricity on a tail.  The proof first builds structure at
selected cardinals, called working levels, and only at the end transfers
it to every larger cardinal.  A tuple is bounded at a level $\mu$ when
its length is strictly less than $\mu$.

\paragraph{From orders to comparison of tuples.}
An order presentation $E$ assigns a model to a linear order: each element
is represented by a label and finitely many indices, and an order embedding
induces a strong map by substituting indices.  Section~\ref{found:section-reversal}
uses specially chosen orders to control the automorphism orbits of short
tuples.  A small subpresentation can be aligned with a larger presentation
so that all its short tuples have matching orbit records.  An elimination
chain then proves that a strong map can be reversed on any prescribed
short tuple.  The small model supplies the alignment; the larger model
accommodates the enumeration of the small model's tuples.  This avoids an
exponentiation hypothesis at the categoricity cardinal.

In Section~\ref{lev:section-diagrams}, the resulting equivalence classes
of tuples under strong maps stabilize as the categoricity cardinal
increases.  These are
the comparison diagrams used later.  Finite supports give two further
tools: a bound on the parameter sets needed to distinguish comparisons,
and simultaneous placement of coherent bounded prescriptions.  Together
they yield amalgamation and homogeneity for bounded tuples at suitable
countably cofinal closure levels.

\paragraph{A calculus that can be applied to systems.}
The same tools will be needed for finite systems of models, not only for
individual models.  Section~\ref{lev:section-working-schemes} records their
precise hypotheses as a working scheme: a finite system shape, specified
arrows, one finite-support presentation, and compatible tuple comparisons.
Sections~\ref{calc:section-calculus} and~\ref{iso:section-isolation} derive
the calculus from this package.  A type over a base is determined by
identification of tuples in an amalgam over that base.  Its distinguished
free extensions are computed by long approximation sequences already in
the base.  The isolation construction gives extensions controlled by
a specified mark: the full type of each additional bounded tuple, together
with the mark, is determined by the mark's type and a bounded fragment of
the base.  Such extensions have the primeness and uniqueness needed for
iteration.

\paragraph{Marked systems detect growth of a geometry.}
Section~\ref{mark:section-increments} shows that adjoining a finite mark
and an extension controlled in this way again gives a working scheme.
Iterating produces a model at each subset of a finite axis set, with
commuting strong maps.  Its boundary consists of the proper-subset
vertices and their maps.  Section~\ref{iter:section-boundaries} first
completes a compatible boundary map into a common extension, then extends
a boundary isomorphism onto the prescribed target.  The return-tree
argument keeps full maps coherent along individual branches but retains
only boundary maps in the common accumulator.  This distinction lets the
orbit bound control joint comparisons across different faces, including
their actual equalities.

Section~\ref{geom:section-geometry} chooses a minimal singleton type,
one with a unique nonalgebraic extension over each relevant base.  Its
realizations form the points of a pregeometry, with independence defined
by finite products of the type.  We retain a distinguished strong copy of
its base model $M_0$, with a fixed diagram for its enumeration, in the
later models; these are the pointed models used below.  Exact lifting for all finite marked
systems at one common level rules out a proper strong extension with no
new points.  A downward construction then extends this no-growth
conclusion to arbitrary larger cardinalities.  Thus the geometry detects
every proper extension, and its basis cardinality equals model
cardinality above the resulting threshold.

\paragraph{From geometric systems to a categorical tail.}
Sections~\ref{sys:section-geometric-systems} and~\ref{exact:section-boundaries}
organize bases into finitely many pieces assigned to vertices.  Tiled
order presentations give systems with these pieces; a paired induction
establishes their working packages and classification by prescribed
basis maps.  A second boundary argument gives exact lifting for every
finite geometric shape at suitable levels.

The two finite-dimensional applications require separate common-level
choices.  Section~\ref{close:section-uniform} supplies both: fixed set codes
for all finite constructions and their bound requests produce a level
where every finite depth is available.  The geometric application is made
after the marked application has supplied the no-growth threshold.
Finally, Section~\ref{tail:section-transfer} starts from a common geometric
level and inducts on the maximum cardinality of the basis pieces,
simultaneously for all finite shapes.  One additional axis records the
previously constructed isomorphism at each extension step; coherent
unions handle limits. The zero-axis case gives pointed categoricity.
The L\"owenheim--Skolem axiom and categoricity at one fixed smaller
cardinal place the chosen base pattern in every sufficiently large model.
Forgetting that base gives the qualitative theorem, and the set-code
reduction gives the uniform theorem.

Figure~\ref{fig:proof-route} summarizes the proof and the two uses of
common-level closure.

\begin{figure}[tbp]
\centering
\begin{tikzpicture}[
  node distance=4mm and 8mm,
  box/.style={draw=linkblue!70,rounded corners=2pt,align=center,
    text width=5.4cm,minimum height=10mm,inner sep=5pt,font=\small},
  closure/.style={box,dashed,text width=4.2cm},
  arr/.style={-{Stealth[length=2mm]},draw=linkblue!80,thick}]
\node[box] (orders) {Unbounded categoricity\\Order presentations and orbit reversal};
\node[box,below=of orders] (diagrams) {Single-model diagrams\\Uniform testing and compact placement};
\node[box,below=of diagrams] (calculus) {Working schemes\\Free extension and relative isolation\\Extensions controlled by a mark};
\node[box,below=of calculus] (point) {Minimal type\\Fixed pointed models};
\node[box,below=of point] (marked) {Finite marked systems\\Exact boundary lifting\\at fixed depth};
\node[box,below=of marked] (growth) {No proper extension with unchanged points\\Basis size equals model size on a tail};
\node[box,below=of growth] (systems) {Geometric systems and relative bases\\Exact lifting for each finite shape};
\node[box,below=of systems] (tail) {Induction through larger cardinalities\\Categoricity on a tail};
\foreach \a/\b in {orders/diagrams,diagrams/calculus,calculus/point,
 point/marked,growth/systems}
  \draw[arr] (\a) -- (\b);
\node[closure,right=of marked] (markedclosure)
  {Common marked level\\All finite marked depths};
\draw[arr,dashed] (marked.east) -- (markedclosure.west);
\draw[arr,dashed] (markedclosure.south) |- (growth.east);
\node[closure,right=of systems] (geometricclosure)
  {Common geometric level\\All finite shapes\\Above the no-growth threshold};
\draw[arr,dashed] (systems.east) -- (geometricclosure.west);
\draw[arr,dashed] (geometricclosure.south) |- (tail.east);
\end{tikzpicture}
\caption{The qualitative proof, ordered by its principal dependencies.
The two dashed boxes are separate applications of
Section~\ref{close:section-uniform}. The first makes all finite marked
depths available for the no-growth argument over the fixed point. The
second is chosen after that argument
and supplies the starting cardinal for the final transfer. These
classes of common levels need not contain a tail. The small-code
reduction then gives the uniform threshold of
Theorem~\ref{main:uniform}.}
\label{fig:proof-route}
\end{figure}


\subsection{Conventions}

All size variables are cardinals unless explicitly called ordinals.
An assertion about bounded tuples at a level $\mu$ concerns tuples of
length strictly below $\mu$. Eligible working levels have countable
cofinality and are strong limits closed under the displayed bounds.
Whenever smaller adequate levels are needed cofinally below a working
level, that requirement is stated and included in the closure bookkeeping.
Automorphisms of an index order always mean automorphisms of its pure
linear order. Charts, auxiliary hulls, and basis choices are not added to
the structures whose isomorphisms are being compared.
